\documentclass[a4paper]{article}
\usepackage[utf8]{inputenc}
\usepackage[margin=1in, footskip=0.25in]{geometry}

\usepackage{graphicx}
\graphicspath{../}

\usepackage{booktabs}

\usepackage[backend=bibtex, style=numeric, sorting=none]{biblatex}
\bibliography{references}
%\addbibresource{references.bib}

\usepackage[font=footnotesize,labelfont=bf]{caption}

\begin{document}
	
	% ======================= start titlepage ==============================================
	
	\begin{titlepage}
		\begin{center}
			\Large Numerical Methods in Astrophysics \\
			\vspace{1cm}
			\huge{
				Project 2 \\
				\vspace{0.5cm}
				\textbf{Ordinary Differential Equations,}\\
				\textbf{Fourier Transforms and} \\
				\textbf{Spectral Analysis} \\
				\vspace{1cm}
			}
			\Large \emph{Saksham Kaushal}
		\end{center}
	\end{titlepage}
	
	% =========================end titlepage ================================================
	
	\tableofcontents
	\newpage
	
	% ========================= Introduction ================================================
	
	\section{Introduction} \label{Introduction}
	
	% ========================= Part 1 ======================================================
	
	\section{Part 1 - Ordinary Differential Equations} \label{task1}
	
		% ----------------------------Methods -----------------------------------------------
		
		\subsection{Methods} \label{methods1}
		
		% ----------------------------Results -----------------------------------------------
		
		\subsection{Results} \label{results1}
		
		\begin{figure} [h]
			\includegraphics[width=\columnwidth]{../task1/q3-epsilon/plots/q3.jpg}
			\caption{Comparison plot for}
			\label{fig:task1.3}
		\end{figure}
		
		\begin{table} [h]
			\centering
			\setlength{\tabcolsep}{20pt} 
			\begin{tabular} {c c}
				\toprule
				\textbf{\(\epsilon\)} & \textbf{\(\sigma_{E_j}\)} \\
				\midrule
				\addlinespace
				\(10^{-5}\) & \(7.326 \times 10^{-6}\)
\\ 
				\(10^{-6}\) & \(6.490 \times 10^{-7}\) \\
				\(10^{-7}\) & \(1.887 \times 10^{-7}\) \\
				\(10^{-8}\) & \(6.918 \times 10^{-7}\) \\
				\(10^{-9}\) & \(1.013 \times 10^{-5}\) \\
				\bottomrule
			\end{tabular}
			\caption{}
			\label{table:task1.1}
		\end{table}
	
		\begin{figure} [h]
			\includegraphics[width=\columnwidth]{../task1/q4-trajectory/plots/q4.jpg}
			\caption{}
			\label{fig:task1.4}
		\end{figure}
	
		\begin{figure} [h]
			\includegraphics[width=\columnwidth]{../task1/q5q6-velocity/plots/q5_combined.jpg}
			\hrule
			\includegraphics[width=\columnwidth]{../task1/q5q6-velocity/plots/q5_sep.jpg}
			\caption{}
			\label{fig:task1.5}
		\end{figure}
		
		\begin{figure} [h]
			\includegraphics[width=\columnwidth]{../task1/q5q6-velocity/plots/q6_combined.jpg}
			\hrule
			\includegraphics[width=\columnwidth]{../task1/q5q6-velocity/plots/q6_sep.jpg}
			\caption{}
			\label{fig:task1.6}
		\end{figure}
		
	% ========================= Part 1 ======================================================
	
	\section{Part 2 - Fourier Transform and Spectral Analysis} \label{task2}
	
		% ----------------------------Methods -----------------------------------------------
	
		\subsection{Methods} \label{methods2}
	
		% ----------------------------Results -----------------------------------------------
		
		\subsection{Results} \label{results2}
		
		\begin{figure} [h]
			\includegraphics[width=\columnwidth]{../task2/plots/q1-5_lin.jpg}
			\caption{}
			\label{fig:task2_lin}
		\end{figure}
		
		\begin{figure} [h]
			\includegraphics[width=\columnwidth]{../task2/plots/q1-5_log.jpg}
		\end{figure}
	
	
	% ========================= Conclusions =================================================
	
	\section{Conclusions} \label{conclusions}
	
	
\end{document}